% GENERATED FILE. Edit canonical lessons, then run npm run generate:books.
% canonical-source-hash: fnv1a64:7693c124913dfb3d
% canonical-lessons: RU-C393-zayavlenie, RU-C393-podavat, RU-C393-notarius, RU-C393-doverennost, RU-C393-zaveshchanie

\chapter{Applications and wills}
\label{ch:ru-applications-and-wills}
\hlchaptermodality{393}

\begin{chapteropening}
\textbf{By the end of this chapter:} \emph{I can put in an application, submit documents, go to a notary, and talk about a power of attorney and a will.}

Complete the last lesson of chapter 393: I can put in an application, submit documents, go to a notary, and talk about a power of attorney and a will.
\end{chapteropening}

\section[\texorpdfstring{\ru{заявление} (zayavléniye)}{заявление (zayavléniye)}]{\ru{заявление} (zayavléniye) — an application (in writing); a statement}

\label{lesson:RU-C393-zayavlenie}



\begin{quote}
Before the new one: say the Russian for an offender, then the Russian for innocent.
\end{quote}

\subsection*{You'll want to know: \ru{заявление}}
\textbf{\ru{заявление}} (\emph{zayavléniye}) — \textquotedblleft{}an application (in writing); a statement\textquotedblright{}.

\begin{quote}
\ru{Он} \ru{подал} \ru{заявление} \ru{на} \ru{отпуск}.
\end{quote}

\begin{quote}
\emph{On pódal zayavléniye na ótpusk.}
\end{quote}

\begin{quote}
He put in a request for leave.
\end{quote}

Neuter. A written request to an office or a boss: \emph{\ru{заявление} \ru{на} \ru{отпуск}}. Also a public statement.

\begin{cousinweb}
\textbf{\ru{заявление}} is built on \emph{\ru{заявить}} \textquotedblleft{}to declare\textquotedblright{}, from \emph{\ru{за}-} and \emph{\ru{явить}} \textquotedblleft{}to show\textquotedblright{}, from Proto-Slavic \emph{*aviti} \textquotedblleft{}to make visible\textquotedblright{} (EtymDB).
\end{cousinweb}

\begin{grammarlens}[title={What you've built}]
One more word for telling, explaining and describing at B1.
\end{grammarlens}

\subsection*{Your turn}
\begin{itemize}
\raggedright
  \item \emph{Say it:} \emph{zayavléniye}
  \item \emph{Say it:} \emph{zayavléniye}, once more
  \item \emph{Say it:} a sentence of your own with \emph{zayavléniye}
  \item \emph{Recall:} say the Russian for an offender, then the Russian for innocent, then say \emph{zayavléniye} again
\end{itemize}

\begin{tcolorbox}[breakable,colback=teal!4,colframe=teal!35!black,title={Before you move on}]
What does \ru{заявление} mean? (\textquotedblleft{}An application (in writing); a statement\textquotedblright{}.) Say it once more.
\end{tcolorbox}
\section[\texorpdfstring{\ru{подавать} (podavát)}{подавать (podavát)}]{\ru{подавать} (podavát) — to submit, to hand in; to serve}

\label{lesson:RU-C393-podavat}



\begin{quote}
Before the new one: say the Russian for innocent, then the Russian for an application.
\end{quote}

\subsection*{You'll want to know: \ru{подавать}}
\textbf{\ru{подавать}} (\emph{podavát}) — \textquotedblleft{}to submit, to hand in; to serve\textquotedblright{}.

\begin{quote}
\ru{Документы} \ru{подают} \ru{до} \ru{первого} \ru{июля}.
\end{quote}

\begin{quote}
\emph{Dokuménty podayút do pérvogo iyúlya.}
\end{quote}

\begin{quote}
Documents must be submitted by the first of July.
\end{quote}

Imperfective; present \emph{\ru{подаю}, \ru{подаёшь}}; perfective \emph{\ru{подать}}. \emph{\ru{Подавать} \ru{заявление}} — to put in an application. \emph{\ru{Подавать} \ru{в} \ru{суд}} — to take someone to court. It also means to serve food.

\begin{cousinweb}
\textbf{\ru{подавать}} is built from \emph{\ru{под}-} \textquotedblleft{}up to, under\textquotedblright{} and \emph{\ru{давать}} \textquotedblleft{}to give\textquotedblright{}, the ongoing partner of \emph{\ru{дать}}: to hand something up to someone.
\end{cousinweb}

\begin{grammarlens}[title={What you've built}]
One more word for telling, explaining and describing at B1.
\end{grammarlens}

\subsection*{Your turn}
\begin{itemize}
\raggedright
  \item \emph{Say it:} \emph{podavát}
  \item \emph{Say it:} \emph{podavát}, once more
  \item \emph{Say it:} a sentence of your own with \emph{podavát}
  \item \emph{Recall:} say the Russian for innocent, then the Russian for an application, then say \emph{podavát} again
\end{itemize}

\begin{tcolorbox}[breakable,colback=teal!4,colframe=teal!35!black,title={Before you move on}]
What does \ru{подавать} mean? (\textquotedblleft{}To submit, to hand in; to serve\textquotedblright{}.) Say it once more.
\end{tcolorbox}
\section[\texorpdfstring{\ru{нотариус} (notárius)}{нотариус (notárius)}]{\ru{нотариус} (notárius) — a notary}

\label{lesson:RU-C393-notarius}



\begin{quote}
Before the new one: say the Russian for an application, then the Russian for to submit.
\end{quote}

\subsection*{You'll want to know: \ru{нотариус}}
\textbf{\ru{нотариус}} (\emph{notárius}) — \textquotedblleft{}a notary\textquotedblright{}.

\begin{quote}
\ru{Мы} \ru{ходили} \ru{к} \ru{нотариусу} \ru{с} \ru{документами}.
\end{quote}

\begin{quote}
\emph{My khodíli k notáriusu s dokuméntami.}
\end{quote}

\begin{quote}
We went to the notary with our documents.
\end{quote}

Masculine, also of a woman. In Russia a notary certifies documents, wills and powers of attorney.

\begin{cousinweb}
\textbf{\ru{нотариус}} is usually traced to Latin \emph{notārius} \textquotedblleft{}a shorthand clerk, a secretary\textquotedblright{} (Lewis \& Short), built on \emph{nota} \textquotedblleft{}a mark, a sign\textquotedblright{}.
\end{cousinweb}

\begin{grammarlens}[title={What you've built}]
One more word for telling, explaining and describing at B1.
\end{grammarlens}

\subsection*{Your turn}
\begin{itemize}
\raggedright
  \item \emph{Say it:} \emph{notárius}
  \item \emph{Say it:} \emph{notárius}, once more
  \item \emph{Say it:} a sentence of your own with \emph{notárius}
  \item \emph{Recall:} say the Russian for an application, then the Russian for to submit, then say \emph{notárius} again
\end{itemize}

\begin{tcolorbox}[breakable,colback=teal!4,colframe=teal!35!black,title={Before you move on}]
What does \ru{нотариус} mean? (\textquotedblleft{}A notary\textquotedblright{}.) Say it once more.
\end{tcolorbox}
\section[\texorpdfstring{\ru{доверенность} (dovérennost)}{доверенность (dovérennost)}]{\ru{доверенность} (dovérennost) — a power of attorney}

\label{lesson:RU-C393-doverennost}



\begin{quote}
Before the new one: say the Russian for to submit, then the Russian for a notary.
\end{quote}

\subsection*{You'll want to know: \ru{доверенность}}
\textbf{\ru{доверенность}} (\emph{dovérennost}) — \textquotedblleft{}a power of attorney\textquotedblright{}.

\begin{quote}
\ru{По} \ru{доверенности} \ru{сестра} \ru{может} \ru{получить} \ru{мою} \ru{посылку}.
\end{quote}

\begin{quote}
\emph{Po dovérennosti sestrá mózhet poluchít moyú posýlku.}
\end{quote}

\begin{quote}
With a power of attorney my sister can collect my parcel.
\end{quote}

Feminine. A paper that lets someone act for you: \emph{\ru{по} \ru{доверенности}}.

\begin{cousinweb}
\textbf{\ru{доверенность}} is built on \emph{\ru{доверить}} \textquotedblleft{}to entrust\textquotedblright{}, from \emph{\ru{до}-} and \emph{\ru{верить}} \textquotedblleft{}to believe\textquotedblright{}, on \emph{\ru{вера}} \textquotedblleft{}faith\textquotedblright{}, Proto-Slavic \emph{*věra} (EtymDB).
\end{cousinweb}

\begin{grammarlens}[title={What you've built}]
One more word for telling, explaining and describing at B1.
\end{grammarlens}

\subsection*{Your turn}
\begin{itemize}
\raggedright
  \item \emph{Say it:} \emph{dovérennost}
  \item \emph{Say it:} \emph{dovérennost}, once more
  \item \emph{Say it:} a sentence of your own with \emph{dovérennost}
  \item \emph{Recall:} say the Russian for to submit, then the Russian for a notary, then say \emph{dovérennost} again
\end{itemize}

\begin{tcolorbox}[breakable,colback=teal!4,colframe=teal!35!black,title={Before you move on}]
What does \ru{доверенность} mean? (\textquotedblleft{}A power of attorney\textquotedblright{}.) Say it once more.
\end{tcolorbox}
\section[\texorpdfstring{\ru{завещание} (zaveshchániye)}{завещание (zaveshchániye)}]{\ru{завещание} (zaveshchániye) — a will (a testament)}

\label{lesson:RU-C393-zaveshchanie}



\begin{quote}
Before the new one: say the Russian for a notary, then the Russian for a power of attorney.
\end{quote}

\subsection*{You'll want to know: \ru{завещание}}
\textbf{\ru{завещание}} (\emph{zaveshchániye}) — \textquotedblleft{}a will (a testament)\textquotedblright{}.

\begin{quote}
\ru{Дед} \ru{написал} \ru{завещание}.
\end{quote}

\begin{quote}
\emph{Ded napisál zaveshchániye.}
\end{quote}

\begin{quote}
Grandfather made a will.
\end{quote}

Neuter. \emph{\ru{Оставить} \ru{завещание}} — to leave a will.

\begin{cousinweb}
\textbf{\ru{завещание}} is built on \emph{\ru{завещать}} \textquotedblleft{}to bequeath\textquotedblright{}, on \emph{\ru{вещать}} \textquotedblleft{}to proclaim\textquotedblright{}, the Church Slavonic form of the root of \emph{\ru{ответ}}, Old Russian \emph{větŭ} \textquotedblleft{}counsel\textquotedblright{}: a will is a last word spoken for the future.
\end{cousinweb}

\begin{grammarlens}[title={What you've built}]
One more word for telling, explaining and describing at B1.
\end{grammarlens}

\subsection*{Your turn}
\begin{itemize}
\raggedright
  \item \emph{Say it:} \emph{zaveshchániye}
  \item \emph{Say it:} \emph{zaveshchániye}, once more
  \item \emph{Say it:} a sentence of your own with \emph{zaveshchániye}
  \item \emph{Recall:} say the Russian for a notary, then the Russian for a power of attorney, then say \emph{zaveshchániye} again
\end{itemize}

\begin{tcolorbox}[breakable,colback=teal!4,colframe=teal!35!black,title={Before you move on}]
What does \ru{завещание} mean? (\textquotedblleft{}A will (a testament)\textquotedblright{}.) Say it once more.
\end{tcolorbox}
